\subsection{Benchmarking against Reference Data}
\label{section:Against reference data}
To validate the implementation the tutorial by \cite{CheckForCorectness} is used as a reference, as it provides a full computation of an principal component analysis based on a small, deterministic dataset with 10 entries. The output of the incremental principal component analysis is compared step by step against the values reported in this reference. The sliding window is configured to span the entire input dataset $w = 10$ and the embedding dimension is set to $d =2$ corresponding to the two-dimensional input data additionally the time delay is set to $\tau = 0$ since the reference example does not involve any temporal lag.
\begin{figure}[htbp]
  \centering
  \setlength{\arraycolsep}{8pt}
  $\text{Data} = 
  \begin{array}{r|r}
    x & y \\
    \hline
    2.5 & 2.4 \\
    0.5 & 0.7 \\
    2.2 & 2.9 \\
    1.9 & 2.2 \\
    3.1 & 3.0 \\
    2.3 & 2.7 \\
    2.0 & 1.6 \\
    1.0 & 1.1 \\
    1.5 & 1.6 \\
    1.1 & 0.9
  \end{array}
  \qquad\qquad
  \text{DataAdjust} = 
  \begin{array}{r|r}
    x & y \\
    \hline
    .69 & .49 \\
    -1.31 & -1.21 \\
    .39 & .99 \\
    .09 & .29 \\
    1.29 & 1.09 \\
    .49 & .79 \\
    .19 & -.31 \\
    -.81 & -.81 \\
    -.31 & -.31 \\
    -.71 & -1.01
  \end{array}$
  \vspace{0.4cm}
  \caption{Reference dataset: original values (left) and mean-adjusted values (right).}
  \label{fig:pca_example_data}
\end{figure}
Due to the reason that the reference tutorial doesn't use time delay embedding or a sliding window approach, a alternative comparison setup is established for validation purposes. The setup executes a direct, one to one comparison between the individual stages of the implementation and the reference calculation. As a first step the data is centered and the mean is updated accordingly. This produces values that match the reference with only small deviations in the resulting covariance matrix. 
\begin{table}[htbp]
\centering
\begin{tabular}{lrrr}
\toprule
Quantity & Reference & Own Implementation & Absolute Difference \\
\midrule
Mean $x$          & 1.81         & 1.81      & $0$ \\
Mean $y$          & 1.91         & 1.91      & $0$ \\
$C_{11}$          & 0.616555556  & 0.616556  & $4.44 \times 10^{-7}$ \\
$C_{12} = C_{21}$ & 0.615444444  & 0.615444  & $4.44 \times 10^{-7}$ \\
$C_{22}$          & 0.716555556  & 0.716555  & $5.56 \times 10^{-7}$ \\
\bottomrule
\end{tabular}
\caption{Comparison of the mean vector and covariance matrix computed by the proposed implementation against the reference values.}
\label{tab:mean_cov_comparison}
\end{table}
In contrast to the reference tutorial, which calculates the eignedecomposition of the covariance matrix directly, the implementation extracts the the two principal components with subspace iteration, which is an iterative method that converges to the true eigenvector after some number of iteration. Since the used dataset is too small to trigger converging, the subspace iteration is executed repeatedly on the same covariance matrix until the resulting eigenvectors stabilize. A single iteration step per incoming samples as in the actual streaming data use case would produce extremely high deltas in the output. When the subspace iteration stabilizes after 200 iterations the resulting deltas are:  
\begin{table}[htbp]
\centering
\begin{tabular}{lrrr}
\toprule
Quantity & Reference & Own Implementation (sign-corrected) & Absolute Difference \\
\midrule
$u_{1,x}$ & 0.735178656  & 0.735179  & $3.44 \times 10^{-7}$ \\
$u_{1,y}$ & 0.677873399  & 0.677873  & $3.99 \times 10^{-7}$ \\
$u_{2,x}$ & 0.677873399  & 0.677873  & $3.99 \times 10^{-7}$ \\
$u_{2,y}$ & -0.735178656 & -0.735179 & $3.44 \times 10^{-7}$ \\
\bottomrule
\end{tabular}
\caption{Comparison of the converged principal components (eigenvectors), sign-corrected due to the inherent sign ambiguity of eigenvectors.}
\label{tab:eigenvector_comparison}
\end{table}
Finally, the mean adjusted data is projected onto the converged principal component with the largest eigenvalue. When comparing the results from the tutorial to the implementation the following are the results. 
\begin{table}[htbp]
\centering
\begin{tabular}{rrrr}
\toprule
Reference $x$ / $y$ & Own Implementation $x$ / $y$ & Abs. Diff. $x$ & Abs. Diff. $y$ \\
\midrule
-0.827970186 / -0.175115307 & -0.827970 / -0.175115 & $2.0 \times 10^{-7}$ & $1.4 \times 10^{-7}$ \\
 1.777580330 /  0.142857227 &  1.777580 /  0.142857 & $3.3 \times 10^{-7}$ & $2.3 \times 10^{-7}$ \\
-0.992197494 /  0.384374989 & -0.992198 /  0.384375 & $5.1 \times 10^{-7}$ & $1.1 \times 10^{-7}$ \\
-0.274210416 /  0.130417207 & -0.274211 /  0.130417 & $5.8 \times 10^{-7}$ & $2.1 \times 10^{-7}$ \\
-1.675801420 / -0.209498461 & -1.675800 / -0.209498 & $1.4 \times 10^{-6}$ & $4.6 \times 10^{-7}$ \\
-0.912949103 /  0.175282444 & -0.912949 /  0.175283 & $1.0 \times 10^{-7}$ & $5.6 \times 10^{-7}$ \\
 0.099109438 / -0.349824698 &  0.099109 / -0.349825 & $1.4 \times 10^{-7}$ & $3.0 \times 10^{-7}$ \\
 1.144572160 /  0.046417258 &  1.144570 /  0.046417 & $1.6 \times 10^{-6}$ & $1.4 \times 10^{-7}$ \\
 0.438046137 /  0.017764630 &  0.438046 /  0.017765 & $1.4 \times 10^{-7}$ & $7.0 \times 10^{-8}$ \\
 1.223820560 / -0.162675287 &  1.223820 / -0.162675 & $5.6 \times 10^{-7}$ & $2.9 \times 10^{-7}$ \\
\bottomrule
\end{tabular}
\caption{Comparison of the projected principal component scores between the reference and the proposed implementation, sign-corrected due to the inherent sign ambiguity of the underlying eigenvectors.}
\label{tab:scores_comparison}
\end{table}

// dazu schreiben ,dass man im grunde keine genaueren werte mit longs bekommen kann zwischen 10 hoch -5 und 10hoch -9 ist aktzeptabel
//vis datensätzte und https://github.com/OliverHennhoefer/awesome-ts-anomaly-detection-datasets
